\subsection{Explicit moduli for the trained neural representation}\label{sec:neural-moduli57}
The transfer bound need not leave acquisition error as an unspecified constant. In the original investment family, write the learned continuation as
\[
 h(y)=v^\top y+c+\sum_{j=1}^{m}u_j\sigma(w_j^\top y+b_j).
\]
Let $s_j=w_j^\top\gamma$, let $R_j$ be the feature's scalar-innovation radius from Lemma~\ref{lem:integrated56}, and set $D_C=\operatorname{diam}G_C$ and $r_C=\sup_{x\in C}\|x-x_C\|_\infty$. Define
\[
 d_j(D)=
 \begin{cases}
 0,&s_j=0,\\
 \min\{1,|s_j|D/(2R_j)\},&s_j\ne0,\ R_j>0,\\
 1,&s_j\ne0,\ R_j=0.
 \end{cases}
\]
All these quantities are determined by the stored coefficients and the declared primitive law.

\begin{proposition}[Action-sensitive neural moduli]\label{prop:moduli57}
For the original separable current action cost and drift with maximum-norm Lipschitz constant $11/16$, the acquisition allowance
\begin{equation}\label{eq:neural-eta57}
 \eta_C=\beta_t\frac{11}{16}r_C
       \sum_{j=1}^{m}|u_j|\,\|w_j\|_1\,d_j(D_C)
\end{equation}
satisfies \eqref{eq:spatial-contrast57}. In particular, the affine part and every exactly action-null ridge contribute zero to this allowance. For the installed zero reference, a valid continuation-error allowance is
\begin{equation}\label{eq:neural-chi57}
 \chi_C=\beta_tD_C\left\{\frac{960}{91}
               +|v^\top\gamma|+\sum_{j=1}^{m}|u_js_j|\right\}.
\end{equation}
Either allowance may be replaced by the minimum of itself and any independently valid sharper enclosure. Computing the displayed bounds requires $O(md)$ arithmetic, with outward rounding when evaluated numerically.
\end{proposition}
\begin{proof}
Write the integrated ridge as $\Phi_R(z)=\E\sigma(z+RZ)$ for $Z$ uniform on $[-1,1]$. For $R>0$, its derivative is zero, then affine with slope $1/(2R)$, then one. Hence the function
\[
 z\longmapsto\Phi_R(z+s_ja)-\Phi_R(z+s_jb)
\]
is Lipschitz with constant at most $\min\{1,|s_j(a-b)|/(2R)\}$. For $R=0$ the difference of the two hinge derivatives is bounded in absolute value by one almost everywhere, giving the same conclusion with the declared convention. For $s_j=0$ the function is identically zero. The preactivation mean varies by at most $\|w_j\|_1(11/16)\|x-x_C\|_\infty$. The affine continuation contrast is $v^\top\gamma(a-b)$, which does not depend on $x$. The state-only current cost cancels, and the current action-cost difference also does not depend on $x$. Summing the ridge contributions and discounting proves \eqref{eq:neural-eta57}.

Under a common innovation, the zero-policy continuation contrast is bounded by $(1920/91)\|\gamma\|_\infty|a-b|=(960/91)|a-b|$. The one-Lipschitz hinge gives a fitted continuation contrast bound $\{|v^\top\gamma|+\sum_j|u_js_j|\}|a-b|$. The triangle inequality and discount give \eqref{eq:neural-chi57}. Independently valid upper bounds may be minimized because each bounds the same nonnegative quantity. The arithmetic count follows from computing the feature projections and norms.
\end{proof}

These formulas connect feature geometry to acquired-state accuracy without assuming that training discovers a null direction. They complement the confidence-qualified learned-direction result: estimated primitive coefficients must still range over the corresponding confidence set. In the present known-law action-search experiment, the direct signed reference enclosure determines adoption and the full-action lower cover determines the displayed Bellman error. The primitive moduli supply an additional certified account, not a retroactive claim that their conservative sum was the stopping criterion.
